\SourceLeaf{312}%
und die angegebenen Zahlen erfassen die Entwicklung bis einschließlich der Größen achter Ordnung. Damit werden \(t,l,b\) selbst bis einschließlich der Größen neunter Ordnung oder bis zur elften Ordnung ausschließlich bestimmt.

Die Entwicklung von \(\log b\) lässt sich auch auf andere Weise ausführen, nämlich nach Potenzen von \(\beta,\tau\) und \(r\). Setzt man
\[
z=\operatorname{tang}y
\]
so wird
\[
y=z-\tfrac13z^3+\tfrac15z^5-\tfrac17z^7+\tfrac19z^9-\text{usw.}
\]
und hieraus
\[
\log y=\log z-\tfrac13zz+\tfrac{13}{90}z^4-\tfrac{251}{2835}z^6+\tfrac{3551}{56700}z^8-\text{usw.}
\]
Wendet man diese Reihe an auf die Gleichung
\[
\operatorname{tang}\tfrac12b=\frac{\operatorname{tang}\frac12t}{\operatorname{tang}B\mathord{.}\operatorname{tang}T}
\]
nachdem man sie mit den Substitutionen
\[
\begin{aligned}
\operatorname{tang}B\mathord{.}\operatorname{tang}T&=\frac\tau\beta\\
\operatorname{tang}\tfrac12t&=\frac{m\tau}{2\BareRoot{(1-\frac14mm\tau\tau)}}
\end{aligned}
\]
in die folgende Form gebracht hat
\[
\frac{\beta m}{2\BareRoot{(1-\frac14mm\tau\tau)}}=\operatorname{tang}\tfrac12b
\]
so ergibt sich
\[
\begin{aligned}
\log b={}&\log\beta+\log m+\tfrac18mm\tau\tau+\tfrac1{64}m^4\tau^4+\tfrac1{384}m^6\tau^6+\tfrac1{2048}m^8\tau^8+\text{usw.}\\
&-\tfrac1{12}\beta\beta(mm+\tfrac14m^4\tau\tau+\tfrac1{16}m^6\tau^4+\tfrac1{64}m^8\tau^6+\text{usw.})\\
&+\tfrac{13}{1440}\beta^4(m^4+\tfrac12m^6\tau\tau+\tfrac3{16}m^8\tau^4+\text{usw.})\\
&-\tfrac{251}{181440}\beta^6(m^6+\tfrac34m^8\tau\tau+\text{usw.})\\
&+\tfrac{3551}{14515200}\beta^8(m^8+\text{usw.})\\
&+\text{usw.}
\end{aligned}
\]
oder, indem man die Glieder gleicher Ordnung zusammenfasst,
